\documentclass{article}
\usepackage{graphicx} % Required for inserting images

\usepackage{amsmath}
\usepackage{graphicx}
\usepackage{booktabs}
\usepackage{ctex}
\usepackage{hyperref}


\title{几何计算前沿_milestone}
\author{王兴辰 余东格}
\date{June 2026}

\begin{document}

\maketitle


\section{选题}
我们选择了基于八叉树或分形结构的三维生成方法这个任务。其核心思想是将三维形状表示为由粗到细的层级结构：模型首先在低分辨率空间中判断大致形状区域，再逐层展开到更高分辨率，从而避免在大量空白区域中进行无效计算。

围绕这一思路，我们实现并研究了两条互补的技术路线，在结构上保持一致、在叶子层表示与条件化方式上有所不同：
\begin{itemize}
    \item \textbf{方法一（occupancy 直接预测，王兴辰）：} 逐层预测每个 voxel 是否被占据，直接以 occupied voxel 作为输出，最终导出为点云。
    \item \textbf{方法二（VQ-VAE 分形生成，余东格）：} 逐层预测八叉树结构（split），叶子层预测 VQ token，再由冻结的 VQ-VAE 解码为 SDF 并经 marching cubes 得到网格。
\end{itemize}
两种方法都采用“由粗到细、父节点展开为 8 个子节点”的统一框架，便于对比不同叶子表示与条件化策略对生成质量的影响。

% ============================================================
\section{方法一：基于 occupancy 预测的分形生成（王兴辰）}

\subsection{核心方法}
该方法采用由粗到细的层级生成思路。模型从一个低分辨率的完整体素网格开始，对每个候选 voxel 预测其是否被占据；被预测为 occupied 的 voxel 会继续展开为下一层的 8 个子 voxel，直到达到设定的最大深度。主要组成部分包括：

\begin{itemize}
    \item \textbf{层级式体素生成 decoder：} 模型从 \texttt{full\_depth} 开始，逐层扩展到 \texttt{depth\_stop}，实现 coarse-to-fine 的三维结构生成。
    \item \textbf{父节点到子节点的特征扩展：} 每个父 voxel 的特征通过局部 cross-attention 模块扩展为 8 个子 voxel 特征，并利用 octant position embedding 区分 8 个不同方向的子节点。
    \item \textbf{Fourier 位置编码：} 模型将 voxel 的三维坐标映射到 $[-1,1]^3$ 空间，并使用 Fourier features 进行位置编码，从而增强模型对空间结构和高频细节的表达能力。
    \item \textbf{全局条件信息：} 训练阶段使用采样点构造全局 shape context，同时加入随机 latent noise，以增强模型的生成能力。为避免模型过度依赖真实点云条件，代码中还加入了 context dropout。
    \item \textbf{occupancy 监督：} 每一层的候选 voxel 都会被监督为 occupied 或 empty。由于三维空间中 empty voxel 数量远多于 occupied voxel，模型使用 Focal Loss 缓解类别不平衡问题。
\end{itemize}

整体生成流程可概括为：
\[
\text{粗层体素网格}
\rightarrow
\text{occupancy预测}
\rightarrow
\text{occupied节点展开}
\rightarrow
\text{下一层occupancy预测},
\]
该过程不断重复，直到生成最终深度的 occupied voxel。

\subsection{当前实现与进展}
训练阶段中，模型根据输入采样点构造每一层的 occupancy 标签：若某个采样点落入一个 voxel，则该 voxel 被标记为 occupied。训练采用 teacher forcing，只有真实 occupied 的节点会被继续展开到下一层，以使训练更稳定。当前总损失由三部分组成：
\[
\mathcal{L}
=
\lambda_{occ}\mathcal{L}_{occ}
+
\lambda_{alive}\mathcal{L}_{alive}
+
\lambda_{rate}\mathcal{L}_{rate},
\]
其中 $\mathcal{L}_{occ}$ 是 voxel occupancy 的 Focal Loss；$\mathcal{L}_{alive}$ 约束被占据的父节点至少生成一个 occupied 子节点；$\mathcal{L}_{rate}$ 约束预测 occupied 比例，避免过密或过稀疏。生成阶段不再使用真实点云，而是依靠随机 latent noise 构造全局条件，逐层保留和展开节点，最终将 occupied voxel 转换为 voxel center point cloud 并保存为 \texttt{.ply} 文件。

截至目前已完成：父节点到 8 个子节点的局部 cross-attention 特征扩展模块、Fourier positional encoding 与 depth embedding、基于采样点的 occupancy label 构造、随机 latent noise / context dropout / FiLM 条件调制；尚未加入 shape context。当前代码已能完成从训练到生成再到可视化的完整流程，但从生成结果看，形状质量仍有提升空间，点云中仍可能存在噪声点或结构不完整的问题。

\begin{figure}[htbp]
    \centering
    \includegraphics[width=0.6\linewidth]{可视化.png}
    \caption{方法一的生成点云可视化。}
\end{figure}

\subsection{当前问题}
\begin{itemize}
    \item \textbf{训练与生成不一致：} 训练用 teacher forcing 依真实节点展开，生成时必须依赖自身预测，一旦前几层出错会在后续层级累积。
    \item \textbf{多样性仍需提升：} 已加入随机 latent noise，但生成结果的多样性与可控性仍需进一步验证与改进。
\end{itemize}

% ============================================================
\section{方法二：基于 VQ-VAE 的分形生成（余东格）}

\subsection{核心方法}
该方法复用 OctGPT 预训练的冻结 VQ-VAE：分形生成器只负责预测八叉树结构与叶子层 VQ token，再由 VQ-VAE 将 token 解码为 SDF，最后经 marching cubes 得到网格。其核心研究问题是：\emph{单次前向（约 4 次 forward）的由粗到细生成器，能否逼近 OctGPT 这类迭代式自回归模型（每样本约 576 次 forward）的质量}，从而以远低的推理成本获得可比形状。主要组成部分包括：

\begin{itemize}
    \item \textbf{由粗到细的 split 预测：} 从 \texttt{full\_depth} 逐层用 OctFormer 预测每个节点是否分裂，展开至 \texttt{depth\_stop}；与方法一共享“父节点经局部 cross-attention 展开为 8 个子节点”的结构。
    \item \textbf{叶子层 VQ token 与冻结解码器：} 叶子层预测 VQ token，由冻结 VQ-VAE 解码为连续 SDF，再由 marching cubes 提取网格，使叶子表示具备更强的表面细节。
    \item \textbf{条件化：} 原模型缺少条件输入、生成趋同，为此引入类别嵌入与 VAE 隐变量 $z$（由真值形状编码得到，生成时从 $\mathcal{N}(0,I)$ 采样），并在每一层注入，使不同类别可生成不同形状。
    \item \textbf{split 监督：} 与方法一类似，节点是否分裂存在严重类别不平衡，采用 Focal Loss 进行监督。
\end{itemize}

\subsection{当前实现与进展}
训练同样采用 teacher forcing：依真值八叉树逐层展开并监督每个节点的 split 与叶子 token。当前已完成：
\begin{itemize}
    \item 搭通“逐层 split 预测 + 叶子层 VQ token 预测 + VQ-VAE 解码 + marching cubes”的完整训练与生成流程；
    \item 通过加入类别条件与 latent $z$ 修复了原先“所有样本几乎相同”的问题，不同类别可生成不同形状；
    \item 在相同冻结 VQ-VAE 下与 OctGPT 做了初步对比：本方法在推理速度上具有明显优势（前向次数约 4 次 vs.\ 约 576 次），但生成质量仍有差距，正在补充 MMD-CD / COV / 1-NNA 等量化评测。
\end{itemize}

\subsection{实验结果与消融分析}
首先在单个飞机样本上进行过拟合验证。该实验用于确认从 split 预测、VQ token 预测到 VQ-VAE 解码和 OBJ 导出的完整 pipeline 是否可学习。已有单飞机过拟合模型在 epoch 300 达到 split accuracy 0.998、VQ accuracy 0.847、test loss 0.319，并已导出可展示的网格结果。这说明在结构分布足够简单时，模型能够准确学习八叉树结构，整体实现是有效的。

在多类别 im-5 数据集上，我们进一步进行了条件化、latent、模型扩容与迭代式 refinement 的消融实验。主要结果如下：

\begin{table}[htbp]
    \centering
    \begin{tabular}{lcccc}
        \toprule
        实验 & 主要改动 & split acc & 最后一层 split acc & VQ acc \\
        \midrule
        C2 & 类别条件每层注入 & 0.812 & 0.756 & 0.652 \\
        V1 & C2 + latent $z$ & 0.836 & 0.780 & 0.647 \\
        Q2 & 768 维大模型 & 0.840 & 0.788 & 0.646 \\
        Q3 & Q2 + split MaskGIT & 0.836 & 0.776 & 0.644 \\
        Q4 & Q2 + sibling attention & 0.841--0.845 & 0.795--0.800 & 0.645--0.647 \\
        S1 & V1 + leaf VQ MaskGIT & 0.840 & 0.791 & 0.639 \\
        \bottomrule
    \end{tabular}
    \caption{VQ-VAE 分形生成路线在 im-5 多类别数据上的主要消融结果。}
\end{table}

从结果可以看出，类别条件与 latent 能明显缓解无条件模型生成趋同的问题；模型扩容和 sibling attention 可以带来小幅结构提升，其中 Q4 是当前结构指标最好的版本。但 split MaskGIT 没有达到预期，反而出现结构变稀疏、碎片增多的问题。leaf VQ MaskGIT 对 nearest-CD 有一定趋势，但该实验使用的是 384 维配置，而 Q4 是 768 维配置，因此不能作为严格同容量对照。

\subsection{当前问题}
\begin{itemize}
    \item \textbf{单次前向的结构上限：} split 预测准确率在当前数据上基本稳定，单纯扩大模型容量或微调结构难以显著突破，反映出单次前向相对迭代式方法的固有限制。
    \item \textbf{与迭代式方法的质量差距：} 相比 OctGPT 的多步迭代解码，单次/少步前向的生成质量仍偏弱，生成网格中存在的空洞/碎片主要来自 split 结构的漏判，而非 VQ-VAE 解码器本身。
\end{itemize}

% ============================================================
\section{后续计划}
两条路线下一阶段都将围绕生成质量与稳定性改进，并完善量化评测：
\begin{itemize}
    \item 减小 teacher-forcing 训练与 free-running 生成之间的差异，缓解误差累积；
    \item 调整各项 loss 权重，缓解过密、过稀疏或噪声生成问题；
    \item 测试不同的 latent conditioning 方式以提升多样性与可控性；
    \item 补齐 MMD-CD / COV / 1-NNA 等量化指标，并在“速度--质量”维度上系统对比两种方法与 OctGPT。
\end{itemize}
总体而言，本阶段已完成 3D 分形生成模型的基础框架搭建，实现了两条路线从训练、生成到可视化的完整流程；后续工作的重点将放在提升生成形状的稳定性、清晰度与多样性，以及完善量化对比上。


\end{document}
